\documentclass[12pt,a4paper]{ctexart}
\usepackage{amsmath,amssymb}
\usepackage{geometry}
\geometry{left=2.8cm,right=2.8cm,top=2.5cm,bottom=2.5cm}

\begin{document}

\begin{center}
    \zihao{3}\heiti 专业422X黄子隽202544011213作业二
\end{center}

\vspace{1em}

\noindent\textbf{班级：}计科4252

\noindent\textbf{姓名：}黄子隽

\noindent\textbf{学号：}202544011213

\vspace{1em}

\noindent\textbf{一、用等式置换（代入）法证明 $(A-B)\cup B=A\cup B$。}

\vspace{0.5em}

\noindent
证明：因为
\[
A-B=A\cap B^c
\]
所以
\[
(A-B)\cup B=(A\cap B^c)\cup B
\]
由分配律得
\[
(A\cap B^c)\cup B=(A\cup B)\cap(B^c\cup B)
\]
又因为
\[
B^c\cup B=U
\]
所以
\[
(A\cup B)\cap U=A\cup B
\]
因此
\[
(A-B)\cup B=A\cup B
\]
故原式成立。

\vspace{1em}

\noindent\textbf{二、设 $A=\{0,1,2,3\}$，$R$ 是 $A$ 上的关系，且}
\[
R=\{(0,0),(0,3),(2,0),(2,1),(2,3),(3,2)\}
\]
\noindent\textbf{给出 $R$ 的关系矩阵和关系图。}

\vspace{0.5em}

\noindent
（1）关系矩阵

按元素顺序 $0,1,2,3$ 排列，得
\[
M_R=
\begin{pmatrix}
1 & 0 & 0 & 1\\
0 & 0 & 0 & 0\\
1 & 1 & 0 & 1\\
0 & 0 & 1 & 0
\end{pmatrix}
\]

\vspace{0.5em}

\noindent
（2）关系图

顶点为：$0,1,2,3$

有向边为：
\[
0\to0,\quad 0\to3,\quad 2\to0,\quad 2\to1,\quad 2\to3,\quad 3\to2
\]

\vspace{1em}

\noindent\textbf{三、设}
\[
A=\{1,2,3\},\quad
R=\{(x,y)\mid x,y\in A,\ x+2y\le 6\},
\]
\[
S=\{(1,2),(1,3),(2,2)\}
\]

\noindent
求：（1）$R$ 的集合表达式；（2）$R^{-1}$；（3）$\mathrm{dom}\,R,\ \mathrm{ran}\,R,\ \mathrm{fld}\,R$；（4）$R\circ S,\ R^3$。

\vspace{0.5em}

\noindent
解：

\noindent
（1）求 $R$ 的集合表达式

由条件 $x+2y\le 6$，分类讨论：

当 $y=1$ 时，
\[
x+2\le 6 \Rightarrow x\le 4
\]
因为 $x\in\{1,2,3\}$，所以得
\[
(1,1),(2,1),(3,1)
\]

当 $y=2$ 时，
\[
x+4\le 6 \Rightarrow x\le 2
\]
所以得
\[
(1,2),(2,2)
\]

当 $y=3$ 时，
\[
x+6\le 6 \Rightarrow x\le 0
\]
无解。

所以
\[
R=\{(1,1),(2,1),(3,1),(1,2),(2,2)\}
\]

\vspace{0.5em}

\noindent
（2）求 $R^{-1}$

将 $R$ 中每个有序对前后交换，得
\[
R^{-1}=\{(1,1),(1,2),(1,3),(2,1),(2,2)\}
\]

\vspace{0.5em}

\noindent
（3）求 $\mathrm{dom}\,R,\ \mathrm{ran}\,R,\ \mathrm{fld}\,R$

由
\[
R=\{(1,1),(2,1),(3,1),(1,2),(2,2)\}
\]
可得
\[
\mathrm{dom}\,R=\{1,2,3\}
\]
\[
\mathrm{ran}\,R=\{1,2\}
\]
\[
\mathrm{fld}\,R=\mathrm{dom}\,R\cup \mathrm{ran}\,R=\{1,2,3\}
\]

\vspace{0.5em}

\noindent
（4）求 $R\circ S,\ R^3$

由定义
\[
R\circ S=\{(x,z)\mid \exists y,\ (x,y)\in S,\ (y,z)\in R\}
\]

已知
\[
S=\{(1,2),(1,3),(2,2)\}
\]

由 $(1,2)\in S$，且 $R$ 中有 $(2,1),(2,2)$，得
\[
(1,1),(1,2)
\]

由 $(1,3)\in S$，且 $R$ 中有 $(3,1)$，得
\[
(1,1)
\]

由 $(2,2)\in S$，且 $R$ 中有 $(2,1),(2,2)$，得
\[
(2,1),(2,2)
\]

所以
\[
R\circ S=\{(1,1),(1,2),(2,1),(2,2)\}
\]

再求 $R^3$。

先求
\[
R^2=R\circ R
\]
由
\[
R=\{(1,1),(2,1),(3,1),(1,2),(2,2)\}
\]
可得
\[
R^2=\{(1,1),(1,2),(2,1),(2,2),(3,1),(3,2)\}
\]

于是
\[
R^3=R^2\circ R=\{(1,1),(1,2),(2,1),(2,2),(3,1),(3,2)\}
\]

\end{document}